\clearpage
\thispagestyle{plain}
\begin{center}
    \vspace{5.0cm}
    \textbf{\Large Abstract}
\end{center}
\vspace{1em}
Active galactic nuclei (AGN) exhibit stochastic variability over a wide range of timescales across the entire electromagnetic spectrum. These brightness variations carry information about the physical structure of the accretion flow and the coupling between the central ionizing source and the surrounding disk, making time-domain studies of AGN light curves an important tool for studying black hole growth and accretion physics on otherwise unresolved scales. Wide-field surveys such as the Zwicky Transient Facility (ZTF), and the forthcoming Legacy Survey of Space and Time (LSST) at the Vera C. Rubin Observatory, now deliver optical light curves for large numbers of AGN.

However, these light curves are irregularly sampled and can contain substantial gaps, which makes reconstructing the underlying signal an underdetermined problem, better addressed by a probabilistic method that returns a distribution of plausible light curves rather than a single estimate. This thesis investigates the use of neural stochastic flows \parencite{Kiyohara2025SNF}, a recently proposed class of generative models that learn the transition density of a stochastic process directly, for the reconstruction of these optical light curves modelled as a damped random walk. The method is trained and evaluated on simulated light curves with physically motivated parameters, and validated against the analytic transition of the damped random walk, which the learned model should reproduce, and against a Gaussian process with the same kernel, the Bayes-optimal reconstruction baseline.

On low-noise simulations, at the $0.01$ mag photometric precision expected of LSST, the method recovers the analytic transition and reconstructs light curves to within about three per cent of the Bayes-optimal Gaussian process in accuracy, with a calibration that is close to it though not exact, reaching a coverage of $0.61$ where the Gaussian process reaches $0.70$. When the sampling and the heteroscedastic noise are instead matched to those of ZTF, the predictive intervals become systematically overconfident, with a coverage near $0.3$ against a nominal $0.68$, and the accuracy falls to roughly twice the error of the Gaussian process. This is traced to the model treating each observed magnitude as if it were noise-free, with no channel through which the per-epoch uncertainties can enter it, and anchoring too tightly to noisy observations. Applied to real ZTF light curves through a masking experiment, in which withheld observations are predicted and the measurement error of each is folded into the comparison, the reconstruction is close to calibrated on the less noisy bands and retains a deficit on the noisiest. A Gaussian process fitted to each curve, which uses the per-epoch uncertainties directly in both its fit and its predictions, is the more consistent baseline, and both improve on linear interpolation. These results identify the handling of measurement noise as the principal limitation of the approach, and motivate a latent formulation, in which the flow describes an unobserved signal related to the observations through an explicit noise model, as the natural next step.


\clearpage